%! Author = grays
%! Date = 3/30/2026

% Preamble
\documentclass[11pt]{article}

% Packages
\usepackage{amsmath}
\usepackage{booktabs}
\usepackage{caption}
\usepackage{float}
\usepackage{graphicx}
\usepackage{listings}
\usepackage{mhchem}
\usepackage{siunitx}
\usepackage{xcolor}
\usepackage[backend=biber,style=ieee]{biblatex}
\usepackage[colorlinks=true, linkcolor=blue, citecolor=blue, urlcolor=blue]{hyperref}
% Listings settings
\lstset{
    basicstyle=\ttfamily\footnotesize,
    keywordstyle=\color{blue},
    commentstyle=\color{green!50!black},
    stringstyle=\color{orange},
    numbers=left,
    numberstyle=\tiny\color{gray},
    stepnumber=1,
    numbersep=5pt,
    backgroundcolor=\color{gray!10},
    showstringspaces=false,
    breaklines=true,
    frame=single,
    rulecolor=\color{gray!50},
    title=\lstname
}
\sisetup{separate-uncertainty=true}
\DeclareSIUnit\angstrom{\text {Å}}

\addbibresource{references.bib}

\title{Spectroscopy}
\date{April 3, 2026}
\author{
    Patrick Geraghty (300349204) \\
    Evan Champagne (300443446)
}

% Document
\begin{document}
\maketitle

    \section{Introduction}\label{sec:introduction}

        Spectroscopy deals with the analysis of materials or gases by
        measurements of their absorption, reflection, or emission spectra.
        In this experiment, a grating spectrometer apparatus---consisting
        of a light source, collimator, grating, and telescope---is used to
        study the emission spectra of low-pressure gases in which atoms are
        excited by an electrical discharge.

        The primary purpose of this experiment is to explore these atomic and
        molecular emissions.
        Firstly, the grating spectrometer is calibrated using the well-known
        emission doublet of sodium (\ce{Na}), which has two closely spaced
        lines at \SI{5889.95}{\angstrom} and \SI{5895.92}{\angstrom}.
        The fundamental principle governing this is the grating equation
        \begin{equation}
            m \lambda = d (\sin \theta)
            \label{eq:grating}
        \end{equation}
        where \( d \) represents the grating constant, \( m \) is the order of
        diffraction, and \( \theta \) is the diffraction angle for a specific
        wavelength \( \lambda \).
        Next, utilizing the calibrated spectrometer, the emission lines of
        a hydrogen source are measured, alongside those of helium, nitrogen,
        and an unknown gas.
        These measurements are then compared against established spectral
        data, and the unknown gas is identified based on its emission
        characteristics.
        Lastly, the hydrogen spectrum is analyzed to determine the Rydberg
        constant, \( R \).
        This is calculated using the quantized energies of an electron in a
        Coulomb potential, given by
        \begin{equation}
            \frac{1}{\lambda} = R \left( \frac{1}{4} - \frac{1}{n^2} \right)
            \label{eq:rydberg}
        \end{equation}
        where \( n \) is the principal quantum number of the Balmer series
        transitions for each observed spectral line.

    \section{Methods}\label{sec:methods}

        \subsection*{Experimental Setup}
        The experimental setup consists of a collimator, a diffraction
        grating (600 lines/mm), and a telescope mounted on a rotating stage.
        A discharge tube is placed at the focal point of the collimator to
        serve as the light source.
        The light passes through the collimator to create a parallel beam,
        strikes the diffraction grating, and is then observed through the
        rotating telescope, allowing for the measurement of the angles
        corresponding to different spectral lines.

        \subsection*{Data Collection}
        The spectrometer is first calibrated using a sodium lamp, where
        a reference initial angle is recorded with the telescope approximately
        normal to the grating.
        The angles for the first-order diffraction lines are measured, and the
        corresponding wavelengths are calculated using the grating equation.
        The calculated wavelengths are then compared to the known values to
        determine whether the starting angle is accurate or if an adjustment
        is necessary.

        After calibration, the sodium lamp is replaced with discharge tubes
        containing hydrogen, helium, nitrogen, and an unknown gas (referred to
        as Unknown E).
        For each gas, the telescope is swept across the viewing field to
        locate the visible emission lines, and the corresponding angles are
        recorded.

        \subsection*{Data Analysis}
         The raw angular data is processed using a Python
        script~\ref{lst:spectroscopy_analysis} that automates the calibration,
        wavelength calculation, and determination of the Rydberg constant for
        hydrogen.

        The script first calculates the theoretical diffraction angle for the
        sodium reference, \( \theta_{\text{ref}} \), using
        Equation~\ref{eq:grating} rearranged for a first-order (\( m = 1 \))
        diffraction:
        \[
            \theta_{\text{ref}} = \arcsin\left( \frac{\lambda_{\text{Na}}}{d} \right)
        \]
        where \( d \) is the nominal grating constant (calculated from the 600
        lines/mm specification) and \( \lambda_{\text{Na}} \) is the average
        wavelength of the sodium.
        The wavelength \( \lambda_{\text{Na}} \) is taken as an average since
        the two lines are very close together, and the primitive
        spectrometer's resolution may not distinguish them clearly.
        The true zero angle of the spectrometer is then determined by
        offsetting the experimentally measured sodium angle by
        \( \theta_{\text{ref}}\), depending on which side of the central
        maximum the measurement was taken.

        For every subsequent emission line observed in the hydrogen, helium,
        nitrogen, and unknown gas, the absolute angle of diffraction,
        \( \theta \), is found by taking the absolute difference between the
        measured angle and the calculated zero angle.
        The emission wavelength, \( \lambda \), is then computed using the
        calibrated angle:
        \[
            \lambda = d \sin(\theta)
        \]

        The hydrogen spectrum is isolated to calculate the
        experimental Rydberg constant, \( R \).
        The script associates the observed colors of the Balmer series with
        their corresponding principal quantum numbers: Red (\( n=3 \)), Cyan
        (\( n=4 \)), Blue (\( n=5 \)), and Violet (\( n=6 \)).
        Using Equation~\ref{eq:rydberg}, a linear regression is performed,
        plotting \( \frac{1}{\lambda} \) against
        \( \left( \frac{1}{4} - \frac{1}{n^2} \right) \) to extract the value
        of \( R \) from the slope.
        Finally, the script calculates the percent error by comparing this
        experimental slope to the theoretical Rydberg constant,
        \( R_{\text{theoretical}} = \SI{1.097373e7}{\per\meter} \) and
        generates a regression plot of the results.

        \lstinputlisting[language=Python,caption={Python script for analyzing emission spectra and determining experimental Rydberg constants},label={lst:spectroscopy_analysis}]{../../src/spectroscopy_analysis.py}

    \section{Raw Data}\label{sec:data}

        The raw data collected from the spectrometer for each gas is
        summarized in Table~\ref{tab:spectral_observations}.
        Each entry includes the gas type, the base angle at which the gas was
        observed, the color of the spectral line, and the corresponding
        observation angle.

        \begin{table}[H]
            \centering
            \caption{Spectral Emission Observations for Various Gases}
            \label{tab:spectral_observations}
            \begin{tabular}{lccc}
            \hline
            \textbf{Gas} & \textbf{Base Angle ($^\circ$)} & \textbf{Line Color} & \textbf{Observation Angle ($^\circ$)} \\ \hline
            Sodium & 295.0 & Doublet (Na) & 313.6 \\ \hline
            Hydrogen & 295.0 & Violet & 293.0 \\
             & & Blue & 277.5 \\
             & & Cyan & 276.6 \\
             & & Red & 269.0 \\ \hline
            Helium & 295.0 & Purple & 277.0 \\
             & & Cyan & 276.0 \\
             & & Green & 275.3 \\
             & & Yellow & 275.0 \\
             & & Orange & 271.6 \\
             & & Red & 268.6 \\ \hline
            Nitrogen & 295.0 & Distinct & 293.0 \\
             & & Blue lines (max) & 278.6 \\
             & & Blue lines (min) & 276.0 \\ \hline
            Unknown E & 295.0 & Blue 1 \& 2 & 277.6 \\
             & & Blue 3 & 277.3 \\
             & & Blue 4 \& 5 & 277.0 \\
             & & Green & 273.0 \\
             & & Orange & 271.6 \\
             & & Red & 269.0 \\ \hline
            \end{tabular}
        \end{table}

    \section{Analysis and Results}\label{sec:analysis}

        The calibration of the spectrometer using the sodium doublet yielded a
        true zero angle of \SI{292.95 \pm 0.17}{\degree}, compared to the
        recorded starting angle of \SI{295.00 \pm 0.17}{\degree}.
        This resulted in a zero-angle error of \SI{2.05 \pm 0.24}{\degree}.
        The nominal grating constant \( d \), was calculated to be
        \SI{1667.67}{\nano\meter}.
        This calibrated zero angle was applied to all subsequent measurements
        to find the true diffraction angles for the observed spectral lines.

        The observed emission lines for the hydrogen Balmer series are
        detailed in Table~\ref{tab:hydrogen_emission_lines}.
        Applying a linear regression to this data, the experimental Rydberg
        constant was determined to be
        \( R_{\text{experimental}} = (1.17 \pm 0.14) \times 10^7\)
        \si{\per\meter}.
        The linear fit to extract this experimental constant is plotted in
        Figure~\ref{fig:rydberg_fit}.

        \begin{table}[htbp]
            \centering
            \caption{Measured angles and calculated wavelengths for the Hydrogen emission spectrum.}
            \label{tab:hydrogen_emission_lines}
            \begin{tabular}{lcccc}
                \toprule
                Color & $n$ & Measured Angle ($^\circ$) & True Angle ($^\circ$) & Wavelength (nm) \\
                \midrule
                Violet & 6 & $293.00 \pm 0.17$ & $0.05 \pm 0.24$ & $1 \pm 7$ \\
                Blue & 5 & $277.50 \pm 0.17$ & $15.45 \pm 0.24$ & $444 \pm 7$ \\
                Cyan & 4 & $276.60 \pm 0.17$ & $16.35 \pm 0.24$ & $469 \pm 7$ \\
                Red  & 3 & $269.00 \pm 0.17$ & $23.95 \pm 0.24$ & $677 \pm 6$ \\
                \bottomrule
            \end{tabular}
        \end{table}

        \begin{figure}[htbp]
            \centering
            \includegraphics[width=0.9\textwidth]{../../out/rydberg_regression}
            \caption{Linear regression to determine the experimental Rydberg constant from the Hydrogen emission lines.}
            \label{fig:rydberg_fit}
        \end{figure}

        Following the hydrogen analysis, the emission spectra for helium and
        the unknown gas were measured.
        The respective measured angles corrected true diffraction angles, and
        computed wavelengths for these gases are presented in
        Table~\ref{tab:helium_emission_lines} and
        Table~\ref{tab:unknown_emission_lines}.
        The emission line for nitrogen were observed but not precisely
        measured due to their diffuse nature, so they are not included in the
        tables.

        \begin{table}[htbp]
            \centering
            \caption{Measured angles and calculated wavelengths for the Helium emission spectrum.}
            \label{tab:helium_emission_lines}
            \begin{tabular}{lccc}
                \toprule
                Color & Measured Angle ($^\circ$) & True Angle ($^\circ$) & Wavelength (nm) \\
                \midrule
                Purple & $277.00 \pm 0.17$ & $15.95 \pm 0.24$ & $458 \pm 7$ \\
                Cyan   & $276.00 \pm 0.17$ & $16.95 \pm 0.24$ & $486 \pm 7$ \\
                Green  & $275.30 \pm 0.17$ & $17.65 \pm 0.24$ & $505 \pm 7$ \\
                Yellow & $275.00 \pm 0.17$ & $17.95 \pm 0.24$ & $514 \pm 7$ \\
                Orange & $271.60 \pm 0.17$ & $21.35 \pm 0.24$ & $607 \pm 6$ \\
                Red    & $268.60 \pm 0.17$ & $24.35 \pm 0.24$ & $687 \pm 6$ \\
                \bottomrule
            \end{tabular}
        \end{table}

        \begin{table}[htbp]
            \centering
            \caption{Measured angles and calculated wavelengths for Unknown Gas E.}
            \label{tab:unknown_emission_lines}
            \begin{tabular}{lccc}
                \toprule
                Color & Measured Angle ($^\circ$) & True Angle ($^\circ$) & Wavelength (nm) \\
                \midrule
                Blue 1 \& 2 & $277.60 \pm 0.17$ & $15.35 \pm 0.24$ & $441 \pm 7$ \\
                Blue 3      & $277.30 \pm 0.17$ & $15.65 \pm 0.24$ & $450 \pm 7$ \\
                Blue 4 \& 5 & $277.00 \pm 0.17$ & $15.95 \pm 0.24$ & $458 \pm 7$ \\
                Green       & $273.00 \pm 0.17$ & $19.95 \pm 0.24$ & $569 \pm 6$ \\
                Orange      & $271.60 \pm 0.17$ & $21.35 \pm 0.24$ & $607 \pm 6$ \\
                Red         & $269.00 \pm 0.17$ & $23.95 \pm 0.24$ & $677 \pm 6$ \\
                \bottomrule
            \end{tabular}
        \end{table}

    \newpage
    \section{Discussion}\label{sec:discussion}

        The primary objective of this experiment was to empirically determine
        the Rydberg constant through the spectral analysis of hydrogen.
        By plotting the inverse of the calculated wavelengths against the
        \(\left( \frac{1}{4} - \frac{1}{n^2} \right)\) term from the Rydberg
        formula, a clear linear relationship was established.
        The slope of this regression yielded an experimental value of
        \( R_{\text{experimental}} = (1.17 \pm 0.14) \times 10^7\)
        \si{\per\meter}.
        When evaluated against the accepted NIST value of
        \SI{1.097373e7}{\per\meter}~\cite{nist_rydberg}, the experimental
        result carries a percent error of approximately 6.6\%.
        Because the theoretical value falls comfortably within the
        experimental uncertainty bounds, the methodology successfully
        validates the quantized energy level model of the hydrogen atom,
        despite the presence of underlying systematic errors in the specific
        wavelength calculations.

        An analysis of the raw wavelength data reveals a consistent systematic
        shift.
        The calculated wavelengths for the hydrogen Balmer series ---
        \SI{677 \pm 6}{\nano\meter} (Red), \SI{469 \pm 7}{\nano\meter} (Cyan),
        and \SI{444 \pm 7}{\nano\meter} (Blue) --- are uniformly higher than
        the accepted values of \SI{656.28}{\nano\meter} (\( H_{\alpha} \)),
        \SI{486.13}{\nano\meter} (\( H_{\beta} \)), and
        \SI{434.05}{\nano\meter} (\( H_{\gamma} \)),
        respectively~\cite{nist_hydrogen}.
        While these three lines adhered strongly to the linear quantum model
        and were used to calculate the Rydberg constant, the measurement for
        the violet line (\( H_{\delta} \)) was explicitly excluded from the
        final regression.
        Its calculated wavelength exhibited an extreme deviation from the
        expected value.
        Because the other three lines possessed a consistent offset rather
        than a massive divergence, this anomaly is highly unlikely to stem
        from baseline systematic errors.
        Instead, it almost certainly arose a discrete error in the angular
        measurement itself.

        The consistent overestimation seen in the red, cyan, and blue lines\
        strongly suggests an inaccuracy in the grating constant, \( d \),
        which was atken as exactly \SI{1667.67}{\nano\meter} based on the 600
        lines/mm specification.
        Any manufacturing tolerance in the grating would systematically scale
        every subsequent calculation.
        This systemic shift could be worsened by angular misalignment; if the
        plane of the diffraction grating was not perfectly orthogonal to the
        incident collimated beam, the effective grating spacing would be
        altered, introducing a cosine-dependent error across all measurements.

        Beyond systematic biases, the precision of the results was constrained
        by instrumental errors.
        The mechanical resolution of the spectrometer's angular scale dictated
        a baseline uncertainty of \(\pm 0.17^\circ\) for every reading.
        This error was compounded by the physical tuning of the collimator
        slit.
        While a wider slit is necessary to capture enough luminosity to view
        the fainter higher-order lines, it inherently broadens the angular
        width of the spectral peaks, making it difficult to align the viewing
        telescope exactly on the center of maximum intensity.
        Finally, the mathematical analysis relied on the vacuum approximation
        of the grating equation, omitting the refractive index of the ambient
        air in the laboratory.
        While this omission only shifts the expected wavelengths by a small
        fraction, it represents a fundamental limitation in the accuracy of
        the employed mathematical model.

        The analysis of the helium spectrum yielded much less precise results
        overall.
        While a few of the calculated wavelengths fell reasonably close to
        expected values, the majority of the measured helium lines deviated
        significantly from standard NIS tables~\cite{nist_helium}.
        This suggests that the baseline systematic errors were compounded by
        the difficulty of resolving helium's denser array of emission lines,
        potentially leading to the misidentification of adjacent
        spectral peaks or higher-order diffractions.

        Despite the broader inaccuracies encountered with the helium data, the
        emission profile for the unknown gas (E) provided enough relative
        structural data to successfully determine its identity.
        By comparing the specific sequence of the calculated wavelengths ---
        particularly the dense cluster of multiple blue lines alongside the
        distinct green, orange, and red emissions --- to standard reference
        spectra~\cite{nist_krypton}, the sample was confidently identified as
        krypton.

        Beyond the quantitative measurements of these distinct gases, a
        qualitative observation of nitrogen gas highlighted the difference
        between atomic and molecular emissions.
        Unlike the sharp, discrete peaks characteristic of the monatomic
        samples, the nitrogen emission appeared as a broad, continuous band.
        This dense clustering arises directly from its molecular structure.
        As a diatomic molecule, nitrogen possesses additional degrees of
        freedom allowing for rotational and vibrational internal motion.
        The quantum mechanical transitions between these numerous
        rotational-vibrational states produce a spectrum with so many closely
        spaced lines that it visually blurs into a continuous band.

    \section{Conclusion}\label{sec:conclusion}

        This experiment successfully utilized a calibrated grating
        spectrometer to analyze atomic emission spectra and validate the
        quantum mechanical model of the hydrogen atom.
        By measuring the diffraction angles of the hydrogen Balmer series and
        applying a linear regression to the Rydberg formula, the experimental
        Rydberg constant was determined to be
        \( R_{\text{experimental}} = (1.17 \pm 0.14) \times 10^7\).
        This result carries a 6.6\% error compared to the accepted NIST
        standard, confirming the theoretical model while highlighting a
        consistent systematic overestimation of wavelengths.
        This systematic shift was primarily attributed to inaccuracies in the
        nominal grating constant and potential orthogonal misalignment of the
        spectrometer.
        Despite these instrumental limitations, the spectral analysis
        methodology proved robust enough to conclusively identify an unknown
        gas sample as krypton based on its unique atomic fingerprint.
        Furthermore, qualitative observations of nitrogen successfully
        demonstrated the physical distinction between the sharp, discrete
        energy transitions of monatomic gases and the continuous
        rotational-vibrational emission bands characteristic of diatomic
        molecules.

\printbibliography

\end{document}